% Part: lambda-calculus
% Chapter: introduction
% Section: lambda-computable

\documentclass[../../../include/open-logic-section]{subfiles}

\begin{document}

\olfileid{lam}{int}{cmp}
\olsection{توابع \usetoken{S}{lambda definable} محاسبه‌پذیرند}

\begin{thm}
\ollabel{thm:lambda-computable}
اگر تابع جزئی~$f$ به‌وسیلهٔ یک ترم لامبدا !!{lambda defined} باشد،
محاسبه‌پذیر است.
\end{thm}

\begin{proof}
فرض کنید تابع~$f$ به‌وسیلهٔ ترم لامبدای~$X$ !!{lambda defined} باشد.
رویه‌ای غیرصوری برای محاسبهٔ~$f$ توصیف می‌کنیم. با ورودی‌های~$m_0$،
\dots،~$m_{n-1}$، ترم~$X \num m_0 \ldots \num m_{n-1}$ را بنویسید.
درختی بسازید: نخست همهٔ کاهش‌های یک‌مرحله‌ای ترم اصلی را بنویسید؛ در
سطح زیر آن همهٔ کاهش‌های یک‌مرحله‌ای آن‌ها (یعنی کاهش‌های دومرحله‌ای ترم
اصلی) را بنویسید؛ و به همین ترتیب ادامه دهید. اگر روزی به یک عددنما
رسیدید، همان را به‌عنوان پاسخ برگردانید؛ وگرنه تابع در آن ورودی تعریف‌نشده
است.

با توسل به تز چرچ می‌توان نتیجه گرفت که این تابع محاسبه‌پذیر است. راهی
بهتر برای اثبات قضیه آن است که توصیفی بازگشتی از این رویهٔ جست‌وجو بدهیم.
برای مثال، می‌توان دنباله‌ای از توابع و روابط بازگشتی اولیه، مانند
«$\fn{IsASubterm}$»، «$\fn{Substitute}$»،
«$\fn{ReducesToInOneStep}$»، «$\fn{ReductionSequence}$»،
«$\fn{Numeral}$» و جز آن تعریف کرد. آن‌گاه رویهٔ بازگشتی جزئی برای محاسبهٔ
$f(m_0, \dots, m_{n-1})$ در پی دنباله‌ای از کاهش‌های یک‌مرحله‌ای می‌گردد
که با~$X \num{m_0} \dots \num{m_{n-1}}$ آغاز و به یک عددنما ختم شود،
و عدد متناظر با آن عددنما را برمی‌گرداند. جزئیات طولانی و ملال‌آور، اما
از جهات دیگر معمول‌اند.
\end{proof}

\end{document}
